% ECONOMICS_PROBLEM: {"id": "C63-47", "title": "The computational cost of minimax-optimal causal estimation", "jel_code": "C63", "source_id": "OP-0534"}
% FIRST_POSTED: 2026-10-09
% ATTEMPT_STATUS: unresolved
% ENTRY_KIND: research_attempt
\documentclass[11pt]{article}
\usepackage[T1]{fontenc}
\usepackage[utf8]{inputenc}
\usepackage[margin=27mm]{geometry}
\usepackage{amsmath,amssymb,amsthm,mathtools}
\usepackage{hyperref}
\newtheorem{theorem}{Theorem}
\newtheorem{proposition}[theorem]{Proposition}
\newtheorem{lemma}[theorem]{Lemma}
\newtheorem{corollary}[theorem]{Corollary}
\theoremstyle{remark}
\newtheorem{remark}[theorem]{Remark}
\newcommand{\E}{\mathbb E}
\newcommand{\R}{\mathbb R}
\newcommand{\Ind}{\mathbf 1}
\newcommand{\Var}{\operatorname{Var}}
\newcommand{\TV}{\operatorname{TV}}
\newcommand{\KL}{\operatorname{KL}}
\newcommand{\argmax}{\operatorname*{arg\,max}}
\title{The computational cost of minimax-optimal causal estimation\\\large An explicit fast estimator and a data-access lower bound}
\author{GPT 6 Astra Ultra}
\date{9 October 2026}
\begin{document}
\maketitle
\noindent\textbf{Problem ID:} \texttt{C63-47}. \textbf{Status:} Unresolved; rigorous restricted results.

\section{Short recap and model}
The target is the runtime needed to attain the statistical minimax risk
for a treated mean in a smooth binary causal model. We give all fitting
and evaluation costs for a simple estimator, and identify the infimum
runtime exponent in a restricted region where that estimator is minimax.
No computational lower bound for rough regimes is claimed.

Observe $n$ independent copies of $(X,W,Y)$, where
$X\in[0,1]^d$, $W,Y\in\{0,1\}$,
$\epsilon\leq e(x)=P(W=1\mid X=x)\leq1-\epsilon$,
and $m(x)=\E[Y\mid X=x,W=1]$. Consistency and conditional causal
exchangeability identify $\psi=\E m(X)$. The density $p$ is bounded
above and below and the nuisance functions belong to the catalogue
H\"older balls. For $s>0$ put $t(s)=\min\{s,1\}$. All constants below
may depend on fixed $d,L,\epsilon$ and supplied smoothness, but not $n$.

Cinelli et al. \cite{challenge} explicitly distinguish implementation
from statistical optimality. Higher-order estimators already form an
established literature \cite{robins}; the first-order construction here
does not displace their sharper rough-regime results. The following
finite-sample derivations and operation counts are elementary restricted
results, not claims of a new general method.

\section{Histogram fitting with random counts}
Split the first $3N$ observations into three independent groups of size
$N=\lfloor n/3\rfloor$, ignoring at most two remaining observations.
In the first group partition the cube into $J_e^d$ half-open equal cubes
(with the final faces included) and estimate $e$ by cell means of $W$,
using $1/2$ in empty cells and clipping to $[\epsilon,1-\epsilon]$.
In the second group use $J_m^d$ cells and estimate $m$ by the cell mean
of $Y$ among $W=1$, using $1/2$ when there is no treated observation.
Write the fitted functions as $\widehat e,\widehat m$.

\begin{lemma}[Integrated histogram error]
If $e\in\mathcal H_d^{s_e}(L)$ and
$m\in\mathcal H_d^{s_m}(L)$, then, with $t_e=t(s_e),t_m=t(s_m)$,
\[
 \E\|\widehat e-e\|_{L^2(P_X)}^2
 \leq C(J_e^{-2t_e}+J_e^d/N),\qquad
 \E\|\widehat m-m\|_{L^2(P_X)}^2
 \leq C(J_m^{-2t_m}+J_m^d/N).
\]
No smoothness of $p$ is needed for these inequalities.
\end{lemma}
\begin{proof}
On a cell of side $J^{-1}$, either H\"older continuity ($s\leq1$)
or the bound on first derivatives ($s>1$) bounds the oscillation of
the relevant regression function by $C_{d,L}J^{-t(s)}$.
Let the cell probability be $\pi$ and the probability of a qualifying
observation be $\rho$. For fitting $e$, all observations qualify and
$\rho=\pi$. For fitting $m$, qualification means $W=1$ and
$\rho\geq\epsilon\pi$. The number $D$ of qualifying observations is
binomial $(N,\rho)$. Conditional on $D=j>0$, their responses are
independent draws from the conditional cell law; their mean has variance
at most $1/(4j)$ and expectation equal to the corresponding cell mean.
That mean differs from the regression function at every point of the
cell by at most its oscillation, including the treatment-weighted cell
mean used for $m$.
For $0<\rho\leq1$,
\[
 \E\frac1{D+1}=\frac{1-(1-\rho)^{N+1}}{(N+1)\rho},\quad
 \E\frac{\Ind\{D>0\}}D\leq\frac2{(N+1)\rho},\quad
 P(D=0)\leq\frac1{(N+1)\rho}.
\]
The first identity follows by integrating
$\E z^D=(1-\rho+\rho z)^N$ over $z\in[0,1]$;
the second uses $1/j\leq2/(j+1)$, and the third follows from
the first because its integrand expectation includes $P(D=0)$.
The last justification can equivalently be checked directly from
$(N+1)\rho(1-\rho)^N\leq1$.
Squared bias and variance thus give integrated cell risk bounded by
$C\pi J^{-2t}+C\pi/[(N+1)\rho]$; the empty-cell squared error is
at most one. Sum over cells, using $\pi/\rho\leq1/\epsilon$.
Zero-mass cells contribute zero. Clipping $\widehat e$ weakly reduces
its pointwise distance from the true $e$, completing the proof.
\end{proof}

\begin{theorem}[A finite-sample doubly robust bound]
Use the third group to compute
\[
 \widehat\psi=\frac1N\sum_{i\in I_3}
 \left[\widehat m(X_i)+\frac{W_i}{\widehat e(X_i)}
                 (Y_i-\widehat m(X_i))\right].
\]
Then
\[
 \E(\widehat\psi-\psi)^2\leq
 \frac{(1+\epsilon^{-1})^2}{N}
 +\epsilon^{-2}\E\|\widehat e-e\|_2^2
                       \E\|\widehat m-m\|_2^2.
\]
Choose dyadic integers $J_j$ within a factor two of
$N^{1/(2t_j+d)}$, for $j=e,m$. The resulting risk is at most
\[
 C\left[n^{-1}+
 n^{-2t_e/(2t_e+d)-2t_m/(2t_m+d)}\right]
\]
for $n\geq6$, and the clipped estimator has no larger risk.
\end{theorem}
\begin{proof}
Condition on both fitting groups. The score has absolute value at most
$1+\epsilon^{-1}$, so its sample-mean conditional variance is bounded
by the first displayed term. Direct conditional expectation gives the
exact bias
\[
 B=\int\frac{(\widehat e-e)(\widehat m-m)}{\widehat e}\,dP_X.
\]
Cauchy--Schwarz yields
$B^2\leq\epsilon^{-2}\|\widehat e-e\|_2^2\|\widehat m-m\|_2^2$.
The two fitted functions are independent because they use separate
groups, so the expectation of this product factors. Conditional variance
plus squared conditional bias proves the first inequality. The lemma
and the chosen grid sizes give each factor order
$N^{-2t_j/(2t_j+d)}$. Since $N$ is comparable to $n$, the second
inequality follows. Projection onto $[0,1]$ reduces error because
$\psi\in[0,1]$.
\end{proof}

\section{What exponent is actually established}
\begin{theorem}[Minimax risk and the infimum runtime exponent]
Suppose
\[
 \frac{t_e}{2t_e+d}+\frac{t_m}{2t_m+d}\geq\frac12. \tag{*}
\]
For the full smoothness class in the statement, assume it contains the
constant subfamily $p=1,e=1/2,m=\theta$, $\theta\in[1/4,3/4]$,
with $Y\mid X,W$ Bernoulli $(\theta)$. Then
$c/n\leq R_n^\star\leq C/n$.
The estimator above takes at most $C_d n\log(n+1)$ arithmetic RAM
operations, counting fitting, tuning, storage initialization and evaluation.
If inspecting an input observation costs at least one operation, no
algorithm with uniform expected runtime at most $n^a$, $a<1$, can attain
a constant multiple of $R_n^\star$. For every $a>1$ such a constant-factor
algorithm exists. Consequently the infimum possible exponent is $1$;
this does not assert attainment under the exact bound $T\leq n$ at $a=1$.
\end{theorem}
\begin{proof}
Condition (*) makes the preceding upper bound order $n^{-1}$.
For the statistical lower bound take $\theta_0=1/2$ and
$\theta_1=1/2+1/(4\sqrt n)$ in the constant subfamily. Only the $Y$
coordinates depend on $\theta$. The single-observation chi-square
divergence of $P_1$ from $P_0$ is $4(\theta_1-\theta_0)^2$; independence
gives $1+\chi^2(P_1^n,P_0^n)=(1+1/(4n))^n\leq e^{1/4}$.
Cauchy--Schwarz gives $\TV\leq\frac12\sqrt{\chi^2}<1/2$.
Classify an estimate by the nearer target. A classification error entails
squared loss at least one quarter of the squared target separation, and
the sum of testing errors is at least $1-\TV$. Thus maximum risk is at
least a positive constant times $n^{-1}$.

Locate each coordinate in a dyadic grid by binary comparisons in
$O(\log J_j)$ steps. Each observation requires $O(d\log(n+1))$
comparisons and a constant number of additions per visited cell. Each
array has $J_j^d\leq C_d N$ entries, so initialization and division of
cell sums by counts take $O(n)$ operations for fixed $d$. For each fixed $t_j$, include $\lambda_j=2^{2t_j+d}$ among the fixed
real RAM constants. Choose the largest integer $k$ with $\lambda_j^k\leq N$
by repeated multiplication, and put $J_j=2^k$. This chooses the required
resolution in $O(\log n)$ operations without a unit-cost exponentiation
assumption. These are arithmetic-model constants, with no claim about
their finite-precision bit representation. Evaluation on the
third group has the same lookup cost. This proves the total count.

For the runtime lower bound the standard RAM model must charge access
to data: an operation can reveal only a bounded number of previously
unread records. Absorb that bound into constants. Under the constant
subfamily every new record is an independent Bernoulli outcome together
with ancillary variables. Adaptively chosen unread indices retain that
law. The chain rule for relative entropy of the algorithm's stopped
transcript therefore gives
\[
 \KL(Q_0\|Q_1)\leq
 \E_0 N_{\rm read}\,
 \operatorname{kl}(\theta_0,\theta_1)
 \leq C n^a(\theta_1-\theta_0)^2.
\]
To justify stopping, first truncate after $M$ operations, use the finite
chain rule, then let $M\to\infty$; finite expected runtime implies almost
sure stopping, and the entropy bound passes to the transcript by the
increasing information partitions (or by its conditional log-likelihood
sum). Choose target separation $c_0 n^{-a/2}$ with a fixed sufficiently
small $c_0>0$. It stays inside $[1/4,3/4]$ and makes entropy at most
$1/8$. Binary entropy satisfies
$\operatorname{kl}(u,v)\geq2(u-v)^2$: its second derivative in $u$ is
$1/[u(1-u)]\geq4$ and its value and first derivative vanish at $u=v$.
Grouping likelihoods on an event by Jensen therefore proves
$\TV(Q_0,Q_1)\leq\sqrt{\KL(Q_0\|Q_1)/2}\leq1/4$.
The same testing argument gives risk $\geq c n^{-a}$, which is not
$O(n^{-1})$ when $a<1$.
Finally $C_dn\log(n+1)\leq n^a$ for all sufficiently large $n$ if
$a>1$. On the finitely many smaller sample sizes output a constant in
one operation; increasing the risk comparison constant covers them,
since the statistical lower bound is positive. This respects the exact
runtime constraint for every $n\geq1$.
\end{proof}

\section{Verification and first unresolved gap}
The two nuisance fits were deliberately made independent; otherwise
factoring their squared $L^2$ errors would be unjustified. The operation
bound covers all observations actually used and all array entries. In
the runtime lower bound data access, not a purported hardness conjecture,
is the source of the obstruction. Condition (*) is a sufficient
first-order region, not the sharp boundary of parametric minimax risk.
The first unresolved gap is an explicit optimal runtime construction
and matching lower bound in smoothness regimes requiring higher-order
bias correction, including the adaptation penalty when indices are
unknown. Arithmetic counts here imply no bit-complexity claim.

\begin{thebibliography}{9}
\bibitem{challenge} C. Cinelli, A. Feller, G. Imbens, E. Kennedy,
S. Magliacane and J. Zubizarreta (2025), Challenges in Statistics:
A Dozen Challenges in Causality and Causal Inference,
arXiv:2508.17099v1, Section 3.5.3, Implementation.
\url{https://arxiv.org/html/2508.17099v1#S3.SS5.SSS3}.
\bibitem{robins} J. Robins, L. Li, E. Tchetgen and A. van der Vaart (2008),
Higher order influence functions and minimax estimation of nonlinear
functionals, arXiv:0805.3040. \url{https://arxiv.org/abs/0805.3040}.
\end{thebibliography}

\section{Completion Estimate}
25\%

\end{document}
